% Image credits for the photographs and AI illustrations of Book 4
% (University Physics, Year 2). Machine-readable license records live in
% images/book4/CREDITS.md; keep the two files in sync.
\chapter*{Kredit Gambar}

Gambar-gambar dalam buku ini adalah karya asli. Foto-fotonya digunakan
berdasarkan lisensi bebas masing-masing, dengan terima kasih kepada para
pembuatnya:

\begin{itemize}
  \item \emph{Pipa pesat pembangkit listrik tenaga air} (Bab 5):
        Chinchu.c, Wikimedia Commons, CC0.
  \item \emph{Deretan pusaran von Kármán pada awan} (Bab 4): Robert
        Cahalan, NASA/GSFC, domain publik.
  \item \emph{Kabel koaksial} (Bab 8): Mx.~Granger, Wikimedia
        Commons, CC0.
  \item \emph{James Clerk Maxwell} (Bab 11): Smithsonian Institution,
        domain publik.
  \item \emph{Aurora selatan dari Stasiun Luar Angkasa
        Internasional} (Bab 14): Joe Acaba, NASA, domain publik.
  \item \emph{Pola interferensi antara dua permukaan datar}
        (Bab 19): majalah \emph{Machinery}, sekitar 1915, domain
        publik.
  \item \emph{LIGO Hanford, tampak dari udara} (Bab 20): LIGO
        Laboratory, domain publik.
  \item \emph{Laser helium--neon pada 594~nm} (Bab 23): Telementor,
        Wikimedia Commons, CC~BY~4.0.
  \item \emph{Matahari dalam ultraungu ekstrem} (Bab 26): NASA/SDO
        (AIA), domain publik.
  \item \emph{Bumi dari Apollo 17} (Bab 26): awak Apollo 17, NASA,
        domain publik.
  \item \emph{Ludwig Boltzmann} (Bab 29): fotografer tak dikenal,
        domain publik.
  \item \emph{Pola difraksi elektron} (Bab 30): National Institute of
        Standards and Technology, domain publik.
  \item \emph{Permukaan emas (100) dilihat dengan mikroskop
        penerowongan payar} (Bab 31): Erwin Rossen, Wikimedia
        Commons, domain publik.
  \item \emph{Charles Townes dan maser pertama} (Bab 31): Dan Rubin,
        domain publik.
\end{itemize}

Tautan sumber dan lisensi tiap foto ada dalam
\texttt{images/book4/CREDITS.md} di repositori buku ini.

\bigskip
Selain foto-foto tersebut, tiga puluh gambar pada Bab 1--10 dan
12--31 adalah ilustrasi buatan kecerdasan buatan, dihasilkan dengan
pembangkit gambar OpenAI dari perintah yang ditulis oleh para
penyunting buku ini, dan diperiksa ketepatan fisikanya sebelum
dimuat. Perintah lengkap untuk setiap gambar tercatat dalam
\texttt{images/book4/ai/PROMPTS.md} di repositori.
\clearpage
